\section{Recovering the ideal outputs and composing extractions}
\label{sec:recovery}

The assignment procedure in Section~\ref{sec:joint} leaves a direct sum
of variable restrictions of the intended outputs.  We now select
sufficiently many of these outputs and recover them exactly.  Two different probability spaces occur: permutations
of an ideal output control inherited masks, whereas the random hash
controls assignment collisions.  Keeping these spaces separate is
essential when a received population already has type restrictions.
Recovery from deletions on all three sides appears in
\cite[Section~7, especially Property~7.1 and Theorem~7.2]{VXXZ2023};
its interface-tensor consequence is Corollary~4.2 there.  The
constituent-stage analysis in \cite[Section~6.5]{MoreAsymmetry2026}
also accounts for inherited approximate type restrictions.  Here we
establish uniform loss on every full-variable orbit of each selected
heterogeneous output and reconstruct its coefficients by an exact linear
restriction.

\subsection{Uniform control of inherited masks}\label{sec:recovery-masks}

Fix all population masses, rational laws, and positive tolerances before
letting the divisible integer $N$ grow.  Classes of zero mass are omitted.
For a target coarse triple $e$, write $Q_e$ for its ideal child product,
including its prescribed child-side type windows, but before inherited
input masks and global pair-frequency masks.  Keep distinct the
population, its previous history, its exact parent split triple, and its
left or right half.  Each resulting child class $c$ has $n_c=c_cN$
positions, with $c_c>0$ fixed.  At an ordinary extraction the corresponding
whole-strand classes, or the unrestricted child classes of an ordinary
split, are used instead.

Let $G_e$ be the product of the symmetric groups on these classes.
A permutation acts simultaneously on the three sides of a child tensor,
and different child classes, including opposite halves, are permuted
independently.  This action preserves both the tensor factors and the
child-side type windows, so it consists of coefficient-preserving
automorphisms of $Q_e$.  The action is on complete variables: a symbol
records the entire child word, including its weight-one labels.

\begin{lemma}[Uniform loss from inherited masks]
\label{lem:orbit-masks}
The total number of $G_e$-orbits of variables on the three sides is
polynomial in $N$, uniformly over target triples $e$.
Suppose every inherited or global mask is either implied by the exact
coarse choices, or tests an empirical ordered pair of child statistics
(or a fixed function of that pair) against its prescribed mixture law
with a fixed positive tolerance.
Choose the child tolerances sufficiently smaller than those mask
tolerances, with the mixture laws consistent as in
Theorem~\ref{thm:joint-step}.  There is then a constant $C_0$ such that,
in every full-variable orbit on every side, the fraction failing any
of these deterministic masks is at most $C_0/N$.
The constant may depend on the fixed data and tolerances, but not on
$e$, the orbit, or $N$.
\end{lemma}

\begin{proof}
Let $\mathcal A_{c,W}$ be the finite alphabet of complete child variables
on side $W$ in class $c$.  An orbit is specified by the multiplicity of
every symbol in every class.  Its number on side $W$ is therefore at
most
\[
 \prod_c(n_c+1)^{|\mathcal A_{c,W}|}=O(N^{d_W})
\]
for a fixed integer $d_W$.  Child type windows merely select some of
these orbits.  The alphabets include all labels, even when a mask uses
only a much coarser statistic.

Consider one mask with tolerance $\eta>0$, on a population of size
$m=cN$.  Its tested law is an ordered-pair mixture.  Write its exact
parent-triple classes as $a$, with weights $w_a=n_a/m$, and let
$q_{a,L},q_{a,R}$ be the two empirical statistic laws in a fixed orbit.
For the prescribed laws $\nu_{a,L},\nu_{a,R}$, the child windows give
coordinate errors at most $\tau$.  If a coarsening of a child statistic
is needed, choose its original window smaller by the fixed alphabet
size, so this assertion holds for the statistic actually tested.
The elementary inequality
\[
 |q_{a,L}(s)q_{a,R}(t)
       -\nu_{a,L}(s)\nu_{a,R}(t)|\leq 2\tau
\]
shows that the mean mixture differs from
$\sum_a w_a\nu_{a,L}\otimes\nu_{a,R}$ by at most $2\tau$
in each coordinate.  More generally, suppose the mask center differs
from this prescribed mixture by coordinatewise error at most
$\alpha\leq\eta/4$.  Choosing $\tau\leq\eta/8$ then leaves a margin
of at least $\eta/2$ inside the mask's acceptance interval.  For the
consistent rational prescriptions of Theorem~\ref{thm:joint-step},
$\alpha=0$.  This margin holds also for orbits at the boundary of a child window.

Uniformly permute an element of the orbit.  Within one parent-triple
class of size $n$, the two half-statistic words are independently
arranged.  If the frequencies of $s$ and $t$ are $a$ and $b$, their
paired count $K_{s,t}$ satisfies, for $n\geq2$,
\begin{equation}
 \mathbb E K_{s,t}=nab,\qquad
 \operatorname{Var}(K_{s,t})
 =\frac{n^2a(1-a)b(1-b)}{n-1}\leq\frac n8.
 \label{eq:matching-variance}
\end{equation}
Indeed, fix the positions of the $na$ occurrences on the left.
The indicators that the right symbol is $t$ have mean $b$,
variance $b(1-b)$, and covariance $-b(1-b)/(n-1)$ at distinct
positions.  Summing gives~\eqref{eq:matching-variance}.  This calculation
is uniform over $a,b\in[0,1]$, including deterministic symbols.

Different parent-triple classes use independent permutations.
Consequently the variance of a pair count summed over a population
is at most $m/8$.  Chebyshev's inequality and the acceptance margin give
\[
 \Pr\left(\left|\frac{\sum_a K_{s,t}^{(a)}}m
       -\mathbb E\frac{\sum_a K_{s,t}^{(a)}}m\right|>\eta/2\right)
 \leq \frac{1}{2c\eta^2N}.
\]
For a mask on a fixed function of the ordered pair, a coordinate of
the pushed-forward law is a sum of at most $B$ pair coordinates,
where $B$ is fixed.  Choose $\tau\leq\eta/(8B)$ and apply the last
bound to each pair coordinate with deviation threshold $\eta/(2B)$.
Summing coordinate errors and taking a finite union bound proves the
same $O(1/N)$ conclusion for that mask.  All positive classes eventually
have size at least two.  There are only finitely many masks and
statistic coordinates, so another union bound gives $C_0/N$.
A uniform group element sends any fixed variable uniformly
over its full orbit; thus this probability is exactly the rejected
fraction of that orbit.  We have not conditioned on passing a mask.
Taking a minimum of the finitely many positive tolerance gaps and
positive population masses makes $C_0$ uniform as asserted.  A statistic
that is itself an ordered pair is simply a symbol of a larger fixed
alphabet.  Masks implied by the exact coarse choices cause no loss.
\end{proof}

\subsection{Exact repair}\label{sec:recovery-exact}

We next give an algebraic recovery statement.  A variable restriction
of a tensor means that some basis variables are set to zero; all
coefficients on the remaining variables retain their original values.

\begin{lemma}[Exact repair by presence masks]
\label{lem:hole-repair}
Let $Q_N$ be a tensor over an arbitrary field, with at most $\exp(bN)$
nonzero monomials for a fixed $b>0$.  Let $G_N$ be a finite group of
coefficient-preserving permutations of its variables, acting separately
on the three sides.  Suppose a variable restriction $H_N$ of $Q_N$
omits at most $C/N$ of each $G_N$-orbit of variables, for a fixed $C>0$.
For all sufficiently large $N$, a direct sum of
\[
 R_N=\exp\bigl(O(N/\log N)\bigr)
\]
copies of $H_N$ has $Q_N$ as a linear restriction.  Every coefficient
is reconstructed exactly once.  The bound is uniform over all
$H_N,Q_N,G_N$ satisfying these same constants.
\end{lemma}

\begin{proof}
We first cover all ideal monomials by shifted copies.  Partitioning
variables by their presence masks then converts this cover into an exact
linear restriction.
Put $\gamma=C/N$ and, for $N>3C$, choose
\[
 L=\left\lceil\frac{(b+1)N}{\log(N/(3C))}\right\rceil
   =O(N/\log N).
\]
Take $L$ independent uniform elements $g_1,\ldots,g_L$ of $G_N$.
For a fixed ideal variable, its probability of absence from $g_iH_N$
is at most $\gamma$.  For a fixed nonzero monomial $xyz$ of $Q_N$,
the union bound on its three variables makes its probability of absence
at most $3\gamma$.  Independence of the $L$ shifts, followed by a
union bound over the support, bounds the probability of an uncovered
monomial by
\[
 \exp(bN)(3\gamma)^L\leq e^{-N}<1.
\]
Fix a family covering every monomial.  The group need only preserve
$Q_N$; it is allowed to move the holes of $H_N$.

For an ideal variable $x$, define its presence mask
\[
 A(x)=\{i\in\{1,\ldots,L\}:x\text{ is retained in }g_iH_N\},
\]
and define $B(y)$ and $C(z)$ similarly.  Partition the three sides
into the classes $X_A,Y_B,Z_C$ of equal masks.  There are at most
$2^L$ classes on each side.  For each triple of classes whose rectangle
$X_A\times Y_B\times Z_C$ contains a nonzero monomial of $Q_N$,
coverage implies $A\cap B\cap C\ne\varnothing$.  Choose one index
$i(A,B,C)$ in this intersection.  All variables in these three classes
are present in $g_{i(A,B,C)}H_N$.  Since this is a variable restriction
of $Q_N$, its restriction to these classes is the full ideal tensor
on the rectangle, with every original coefficient.

Use a separate copy of $H_N$ for each such rectangle, permute that
copy by the chosen $g_{i(A,B,C)}$, and zero all other variables.  In
the direct sum of these rectangular tensors, map the copied variable
$(A,B,C,x)$ to the ideal variable $x$, and similarly on the other two
sides.  These are ordinary linear maps.  Every nonzero monomial $xyz$
belongs to exactly one triple of mask classes, so its coefficient is
supplied once and only once.  Identifying a variable used in several
rectangles introduces no additional monomials: linear maps act on
the sum of the existing tensor terms.  There is no division or
characteristic restriction.  At most $2^{3L}$ copies are used, and
padding with unused copies gives the asserted common bound $R_N$.
\end{proof}

\subsection{Completing the joint extraction}\label{sec:recovery-completion}

\begin{proof}[Completion of the proof of Theorem~\ref{thm:joint-step}]
Use the hash slack $\delta>0$ from Section~\ref{sec:joint}, with all
entropy errors chosen smaller than a fixed fraction of this slack.
For a target triple $e$, let $S_e$ be its own hash-survival event;
$\Pr(S_e)=h_M=|U|/M^2$.  Let $P_e(v)$ denote the deterministic
assertion that an ideal variable $v$ of $Q_e$ passes its inherited and
global masks.  Membership in $Q_e$, together with the tolerance
inclusions, supplies usefulness and compatibility with its own triple.
For every $G_e$-orbit $\mathcal O$, Lemma~\ref{lem:orbit-masks} gives
\[
 d_{e,\mathcal O}
 =\frac{|\{v\in\mathcal O:\neg P_e(v)\}|}{|\mathcal O|}
 \leq C_0/N.
\]

For a passing variable $v$, use the event $R_{e,v}$ defined at the
end of Section~\ref{sec:joint}: it includes failure of $e$ at the $X$
uniqueness test and all compatible-target collisions relevant to $v$.
Equation~\eqref{eq:joint-conditional-loss} gives
\begin{equation}
 \Pr(R_{e,v}\mid S_e)\leq e^{-cN}
 \label{eq:conditional-variable-loss}
\end{equation}
for one fixed $c>0$, uniformly in $e$ and $v$, for all large $N$.
Here $P_e(v)$ is fixed before choosing the hash; the conditioning is
still only $S_e$.  We now average this individual loss over each full
ideal orbit, including in the denominator the variables that fail the
deterministic masks.

Specifically, put
\[
 r_{e,\mathcal O}
 =\frac{1}{|\mathcal O|}
   \sum_{\substack{v\in\mathcal O\\P_e(v)}}
       \mathbf1_{R_{e,v}}.
\]
By linearity of expectation,
$\mathbb E(r_{e,\mathcal O}\mid S_e)\leq e^{-cN}$.
Variables failing $P_e$ contribute to $d_{e,\mathcal O}$, and need
no collision estimate.  If $P(N)$ bounds the total number of orbits
on all three sides, Markov's inequality and a union bound imply
\[
 \Pr\left(\exists\mathcal O:r_{e,\mathcal O}>1/N\mid S_e\right)
 \leq NP(N)e^{-cN}=e^{-cN+O(\log N)}.
\]
Declare $e$ good if it survives hashing, remains eligible, and all
these random fractions are at most $1/N$.  Including the exponentially
small eligibility failure explicitly changes only the polynomial
factor in this bound.  Every good $e$ therefore has missing fraction
at most
\begin{equation}
 \gamma_N=(C_0+1)/N
 \label{eq:uniform-good-holes}
\end{equation}
in each orbit.  The eligible list remains the fixed list used in
Section~\ref{sec:joint}.  Deterministic masks are accounted for as
additional deletion flags; they do not trigger a recomputation of
that list or of the compatibility candidates.

There are $\exp(NA_{\mathrm{amb}}+O(\log N))$ target triples, and
$\log M=N(A_{\mathrm{amb}}-E+\delta)+O(1)$ with
$\log|U|=\log M-o(N)$.  If $\mathcal G$ is the set of good triples,
\begin{align*}
 \mathbb E|\mathcal G|
 &\geq\exp(NA_{\mathrm{amb}}+O(\log N))\,h_M
       \bigl(1-e^{-cN+O(\log N)}\bigr)\\
 &=\exp\bigl(N(E-\delta)-o(N)\bigr).
\end{align*}
Thus one hash choice has at least this many good triples, up to
integer rounding.  This uses no independence between different
targets or between their losses.  The assignment proof has already
made their holed outputs $H_e$ disjoint.  Discard all other outputs.

All strand lengths and alphabets are bounded, and there are $O(N)$
factors, so every ideal $Q_e$ has at most $\exp(bN)$ monomials for
one fixed $b$.  For example, a product on at most $aN$ original
$\CW_5$ sites has at most $18^{aN}$ terms before any masks, since
$\CW_5$ has $3\cdot5+3=18$ terms; any $b>a\log18$ suffices.
Apply Lemma~\ref{lem:hole-repair} with
$C=C_0+1$ to every good output.  A common
$R_N=\exp(O(N/\log N))$ works for all of them.  Indeed,
\[
 \bigoplus_{r=1}^{R_N}\ \bigoplus_{e\in\mathcal G} H_e
 \cong
 \bigoplus_{e\in\mathcal G}\ \bigoplus_{r=1}^{R_N}H_e.
\]
The same replication slots supply each summand with $R_N$ copies.
Within a summand choose its own covering family and rectangular
maps, and keep variables of different summands distinct.  This
produces $\bigoplus_{e\in\mathcal G}Q_e$ with $R_N$ copies of the
original joint input.  The number of summands does not multiply the
replication cost.

Finally, exact target types give the same number of factors in every
population, history, split-triple, and half class for every $e$.
After repair, simultaneous position relabelings within each population
put these classes in a fixed canonical order.  They identify all
$Q_e$ with a single ideal product $Q_N$, without permuting the physical
tensor sides.  This canonicalization concerns the repaired ideal
tensors and their child windows; no invariance of an earlier hole
pattern is needed.  Choose $0<\delta<\min(E,\varepsilon)$ and then $N$ large
enough to absorb the $o(N)$ term.  The result is at least
$\exp(N(E-\varepsilon))$ identical ideal outputs from
$\exp(o(N))$ input copies, as claimed.
\end{proof}

\subsection{Composition and subdivision}\label{sec:recovery-composition}

To iterate Theorem~\ref{thm:joint-step}, we record the two algebraic
operations explicitly.  Write $D_r(T)=\bigoplus_{i=1}^r T$,
and write $A\preceq B$ when $A$ is a linear restriction of $B$.

\begin{lemma}[Replication and composition]
\label{lem:replication}
Suppose $D_{s_1}(T)\preceq D_{r_1}(S)$ and
$D_{s_2}(V)\preceq D_{r_2}(T)$.  Then
\[
 D_{s_1s_2}(V)\preceq D_{r_1r_2}(S).
\]
For any finite collection of tensors,
\[
 \bigotimes_{i=1}^j D_{r_i}(F_i)
 \cong D_{\prod_{i=1}^j r_i}\left(\bigotimes_{i=1}^jF_i\right).
\]
The same assertions permit different maps in different direct-sum
branches, provided a common upper bound on their replication counts
is used.  Consequently every fixed finite sequence of the joint
extractions above, including finite tensor products of them, costs
only $\exp(o(N))$ replications of its original input.
\end{lemma}

\begin{proof}
Take $r_2$ copies of the first restriction and reorder their direct
sum as $D_{s_1}(D_{r_2}(T))$.  Apply the second restriction separately
to its $s_1$ disjoint blocks.  This gives the first identity with
$r_1r_2$ source copies, regardless of the value of $s_1$.
The second identity is direct-sum distributivity: its summands are
indexed by the choices of one replication slot for each factor.
A finished factor can be tensored through with replication count one.
Branch-dependent maps act block by block after padding unused slots
to the common bound.  Lemma~\ref{lem:hole-repair} provides such a
bound uniformly over the selected outputs of each joint step.
Products of finitely many $\exp(o(N))$ quantities are $\exp(o(N))$.
For a fixed finite set of population specifications, all positive
sizes are $cN$ with $c>0$ fixed, so a local $o(cN)$ is also $o(N)$.
\end{proof}

There is also an exact way to subdivide a received population before
its next use.  Let $B$ be its strand tensor, let $\nu_W$ be the
prescribed statistic law on side $W$, and denote by
$U(m,\nu,\eta)$ the restriction of $B^{\otimes m}$ to the three
coordinatewise type windows with tolerance $\eta$.  Partition the
positions into fixed parts of sizes $m_1,\ldots,m_j$, and on each part
use the same laws and a positive tolerance $\eta_i\leq\eta$.  There is an exact
variable restriction
\begin{equation}
 \bigotimes_{i=1}^j U(m_i,\nu,\eta_i)
 \preceq U\left(\sum_{i=1}^j m_i,\nu,\eta\right).
 \label{eq:subdivision-restriction}
\end{equation}
To see this, the empirical law of a concatenated variable is the
weighted average of its partwise empirical laws, with error at most
$\sum_i(m_i/\sum_jm_j)\eta_i\leq\eta$.  Hence every retained
variable satisfies the original window; after the partwise masks are
imposed, no condition couples different parts, giving exactly the
displayed product.  Different side tolerances are handled in the same
way.  Apply this argument separately within every received history
class; classes with different prescriptions remain distinct.

Equation~\eqref{eq:subdivision-restriction} need not retain a large
fraction of the earlier variables.  The next application of
Theorem~\ref{thm:joint-step} starts from this exact product, with its
new positive population masses and its consistent inherited laws.
Thus it supplies its own yield and mask analysis.  Together with
canonicalization after repair, this explains why arbitrary fixed
finite subdivisions and heterogeneous populations can be used in a
subsequent joint step.  All such subdivisions and tolerances are fixed
before $N$ grows; no assertion of uniformity in a number of populations
growing with $N$ is needed.

The required tolerance hierarchy exists for any fixed finite
calculation.  At a step first make its global and compatibility
tolerances small enough for the entropy errors allowed by its hash
slack.  Then choose its child windows smaller than all those tests and
all inherited windows that they must satisfy.  Repeat toward the
descendants.  There are finitely many constraints, each allowing a
strictly positive choice.  All constants in this section may depend
on this fixed calculation; the strand limit is taken only after these
choices.
